\chapter{36 块的产线}
\label{ch:parser-production-line}

第~\ref{ch:branch-pipeline}~章讲的是围栏：分支、PR、回归、止损。这一章讲围栏里到底能生产出什么。答案不是一段演示代码，而是一条能持续迭代、能验证、能交接的解析器生成器产线。

% ---------------------------------------------------------------------
\section{LL(1) 塑料轮子不是妥协，是轨道}

第~\ref{ch:script-engine}~章的 LL(1) 生成器看上去很朴素：文法有限、模式单一、输出的 C 代码也不算复杂。但它的价值从来不只是“能跑一个示例”。它给出了骨架、接口和测试，让后面所有工作都有落点。

有轨道，Agent 才能在你搭的结构上迭代，而不是每次从零猜。词法接口放在哪里，语义动作如何翻译，生成产物怎样验证，哪些能力必须保持一致，这些问题一旦被塑料轮子固定下来，增量扩展就有了参照。先造塑料轮子，不是妥协，是给 Agent 搭第一条产线。

% ---------------------------------------------------------------------
\section{20 块重造 Yacc}

真正的工作从 LL(1) 往后推进。先把直接左递归、预测表和语义动作翻译做稳，再扩展到 SLR(1)，补上 LR(0) 项集和 FOLLOW 集；接着做 LALR(1)，加入 LR(1) 项集与同心合并；最后让优先级、\texttt{\%prec} 和 \texttt{\%expect} 进入同一套冲突消解流程。

这不是一次到位，而是在同一套验证上做增量。每一步都要有 parity check 兜底：新模式的表和结果必须与上一步的已知行为对齐，不能因为加了能力就悄悄改变旧文法。20 块买到的不是一个更大的代码库，而是一条从 LL(1) 通向 LALR(1) 的可验证路线。

% ---------------------------------------------------------------------
\section{5 块重造造 Yacc 的机器}

当功能路线走通之后，问题就从“能不能出结果”变成“机器本身够不够稳”。原来的 Python 原型便于探索，但字符串拼接、容器查找和模板加载方式会在大文法上暴露代价。于是用 C++17 重写，把 Qt 风格命名、Doxygen 注释、模板内联和 StringBuilder 一并整理进去。

这不是换功能，是换机器本身。解析逻辑不变，但代码组织、生成路径、构建方式和运行性能都换了基础。塑料轮子继续负责验证行为，新机器负责把同一套行为跑得更快、更清楚。

% ---------------------------------------------------------------------
\section{36 块把 IELR 和 GLR 也装进去}

最后补上两条硬路线。IELR 位于 LALR 和规范 LR(1) 之间：先构造 LALR 状态，再按冲突需要拆分状态，尽量保留 LALR 的表规模，同时获得规范 LR(1) 的语言覆盖能力。

GLR 则走另一条路：不再要求状态合并后只有一个动作，而是把冲突保留到运行期，用图结构栈并行探索多条路径。它接受更高的时间代价，最坏情况下是 \(O(n^3)\)，换来二义文法、悬空 else 这类场景下的多结果能力。

这两条路线在 Bison 里通常是分开实现的，而这个业余项目在一个生成器里同时做了。36 块，覆盖了解析器生成的主流硬路线。

% ---------------------------------------------------------------------
\section{为什么成本能压到 36 块}

关键不是让 Agent 从零重来，而是让它做增量。LL(1) 的骨架还在，词法接口还在，语义动作翻译还在，分文件布局还在，parity check 也还在。新增模式只替换分析表构造和运行时策略，前端、后端和验证方式尽量复用。

同一套结构、同一套测试、同一套接口，意味着每次新增能力都只改局部。塑料轮子的价值在这里兑现：不是它永远够用，而是它让后面的钢铁轮子不必重新铺轨道。不是 AI 便宜，是你没让它走弯路。

% ---------------------------------------------------------------------
\section{AI 是体力劳工，也是迭代引擎}

Agent 很适合做这种增量：补一个模式，跑一组回归，找出差异，解释失败，再继续补。它不知疲倦，你不知妥协；它负责跑，你负责停。每次它说“已经完成”时，你都要问一句：完成的是当前版本，还是整条产线？

人仍然掌握节奏。什么时候从 Work 切到 Right，什么时候允许进入 Fast，什么时候因为失败次数过多而止损，什么时候冻结接口准备发布，这些都不能外包。它把塑料轮子推到钢铁轮子，你决定什么时候停、停在哪个版本。

% ---------------------------------------------------------------------
\section{轮子不朽}

技术会旧，轮子会被重造。今天用 C++17 写的生成器，过几年也可能被新的语言和工具替代。但“先活下来、再谈优雅”的决策框架不会过时：先做能跑的骨架，再用回归把它变正确，最后在不破坏边界的前提下优化。

STORY.md 里说，轮子归你们了，造轮子的机器也归你们了。真正要留下来的，不只是生成器本身，还有它怎样从塑料轮子长成产线的过程。真正不朽的不是轮子，是造轮子的判断力。

% ---------------------------------------------------------------------
\section{从荒野到产线}

荒野求生，是怕被公司优化掉。建产线，是让公司离不开你的产线。前者是防守，后者是进攻。

之前我的生存方式是影子模式：中午自费、错峰跑 DeepSeek、窗口伪装，不挑战组长的 GPT-6 信仰。那是个体对抗组织低效的无奈之举，是野路子，带不走，也随时可能被掐断。

但产线不一样。它天生就是为公司环境设计的。

老板要的免责，产线里有。AI 干活，人负责，PR 闸门，人工验收。老板一听，安全，放心。组长要的控制权，产线里有。AI 永远不碰 \texttt{main}，只能提 PR，关键用例人工复核。他不仅不反对，还会觉得你做事严谨。公司要的降本增效，产线里有。一个人花 36 块，带着 Agent 把解析器生成器的理论路线全走通了；同样的事，传统团队可能要几个月。这不是个人自费省钱，这是可复制、可量化、能写进季度汇报的工程效能提升。团队要的资产沉淀，产线里有。Qt 框架接口分离、parsergen 的 Git 仓库、\texttt{AGENTS.md} 的规范，全都是能留下的组织资产。你休假了，实习生拉个分支就能接着干，不会因为“只有你知道怎么跑”而卡住。

但是，产线要在公司活下来，还得跨过三道坎。

第一道，不要试图叫醒装睡的人。不要一上来就推销“我的 AI 工作流比 GPT-6 好”。你要包装成：为了保障 AI 生成代码的质量，我建议我们引入工程化的 PR 闸门和全量回归。用公司的痛点去卖你的方案。

第二道，把个人工具变成团队规范。你的 AI 分支、隔离目录、parity check，写进团队的 README 或开发规范里。一旦组长或老板在评审会上说了一句“以后 AI 改代码都走这个流程”，你的产线就正式转正了。

第三道，防白嫖，保护自己。公司可能会享受你的产线带来的效率，但把功劳归给“公司流程好”。你要留好自己的作品集：parsergen、Qt 框架的接口层。这些是你的资产，不是公司的。你可以在公司用，但要把它们整理好，放在自己的 Git 仓库里。

所以，它变成产线以后，你的生存姿态就变了。以前你是一个人偷偷在荒野里求生；现在你是带着一套能复用、能交接、能兜底的产线，去跟公司谈条件。你不再需要自费买 token 躲着用了。你可以大大方方地在工作群里说：“我跑了一轮 AI 全量回归，有问题，AI 已经按 PR 流程打回了，明天人工复审。”

荒野求生，是怕被公司优化掉；建产线，是让公司离不开你的产线。前者是防守，后者是进攻。

% ---------------------------------------------------------------------
\section{项目地址}

项目仓库：\url{https://gitee.com/rabbitc999/Script-Engine}

仓库包含 LL(1)、SLR、LALR、IELR、GLR，支持 \texttt{\%tree}、\texttt{\%ast\_*}、\texttt{\%token regex}、\texttt{\%identifier}、语义动作翻译、分文件布局和 parity check。读者可以克隆仓库，先用 \texttt{example/*.gram} 跑一遍，再对照本章理解每种模式的取舍。

运行方式以仓库 README 为准。最小路径是先构建生成器，再用一个示例文法生成解析器工程：

\begin{lstlisting}
make
./parsergen --grammar-file example/calc.gram --out-dir example/calc_out
\end{lstlisting}

仓库中的 \texttt{RESULT.md} 会记录实跑结果，Makefile 和 CMakeLists.txt 用于复现构建。项目采用 GPL-3.0 发布，使用、修改和再发布时请遵守相应条款。
